\chapter{Abstract} \begin{tabular}{@{} >{\bfseries}p{0.2\textwidth} p{0.77\textwidth} @{}} \MakeUppercase{Keywords} & sepsis-associated acute kidney injury (SA--AKI); sepsis-3; KDIGO; early warning; risk stratification; leakage-aware modeling; calibration; SHAP explanations; decision-curve analysis; survival analysis; MIMIC--IV; UMLS; RxNorm; SNOMED CT; ATC\\ \end{tabular} \par \begin{doublespacing} This thesis addresses a critical, time-sensitive clinical question: at the end of the first ICU day (T0--T24), can we identify which patients with sepsis-associated acute kidney injury (SA--AKI) are at the highest risk of in-hospital death, and can we explain the reasoning? Previous SA--AKI models, while often reporting strong discrimination, suffer from methodological issues that limit their clinical utility, including temporal leakage from using post-baseline data, global fitting of preprocessing steps, and a lack of model calibration. To overcome these limitations, we developed a leakage-averse, TRIPOD-aligned framework. An ontology-integrated data warehouse was built using PostgreSQL to harmonize the MIMIC--IV dataset with medical terminologies (UMLS, RxNorm, SNOMED CT, ATC). The SA--AKI cohort was rigorously defined using Sepsis--3 and KDIGO criteria, with all predictors strictly confined to the T0--T24 window. A one-time 80/20 stratified split ensured that all data-dependent steps—including a novel approach to creating and selecting significant missingness indicators and performing median imputation—were fitted only on training data. The methodology encompassed three stages: (i) an extensive exploratory data analysis on the training set to identify clinical phenotypes; (ii) supervised mortality prediction using a benchmark of models, from which a LightGBM classifier was selected, tuned with Optuna, and calibrated using an isotonic regressor on a validation set; and (iii) survival analysis to model time-to-death using the same early covariates with Cox Proportional Hazards and a custom DeepSurv model. On the held-out test set, the final calibrated LightGBM model achieved an AUROC of \textbf{0.749} and an F1-score of 0.522, demonstrating strong and reliable discriminative ability. The survival analysis yielded C-indexes of 0.69 (Cox) and 0.69 (DeepSurv). The primary contribution of this work is a reproducible, timing-faithful pipeline that delivers a calibrated, interpretable risk score and time-to-event estimates, making it a valuable and trustworthy tool for bedside clinical decision support. \end{doublespacing}